\originalchapter{割、连通度与块}\label{ch:connectivity}
\sourceof{18.315 L23; 块分解与排序细节为编者补充}
\originalsection{边割与顶点割}
连通只说明还能通行, 没有衡量网络对故障的承受能力. 删除一条边或一个顶点可能产生完全不同的结果. 顶点故障连带删除所有关联边, 因而需要区分两种连通度.
\begin{definition}
对 $\varnothing\ne S\subsetneq V$, 边割 $\partial S$ 是一端在 $S$、另一端在补集的边集. 连通图中删去后不再连通的单条边称为桥. 若删去顶点 $v$ 增加分支数, 则 $v$ 是割点. 连通非完全图的顶点连通度 $\kappa(G)$ 是使图不连通所需删除的最少顶点数; 完全图约定 $\kappa(K_n)=n-1$. 边连通度 $\lambda(G)=\min|\partial S|$ ($n\ge2$). 非连通图的两种连通度均记0.
\end{definition}
\begin{proposition}
一条边是桥当且仅当它不在任何圈上. 对 $n\ge2$ 的连通图, $\kappa(G)\le\lambda(G)\le\delta(G)$.
\end{proposition}
\begin{proof}
若边在圈上, 圈的另一段提供替代通路, 故不是桥. 若非桥, 删边后其两端仍由路相连, 加回该边就产生圈. 单顶点割的边数为它的度, 所以 $\lambda\le\delta$.

设最小边割分开非空 $A,B$, 其大小为 $k$. 如果存在 $a\in A,b\in B$ 不相邻, 对每条跨割边选择一个端点删除, 避免选择 $a,b$: 含 $a$ 的边删另一端, 含 $b$ 的边亦然, 其余任选. 至多删 $k$ 个顶点后 $a,b$ 分离. 若两侧完全相连, 则 $k=|A||B|\ge |A|+|B|-1=n-1\ge\kappa$. 得到所需不等式.
\end{proof}
一条桥的两端不一定都是割点: 单边图的任何顶点删除后只剩单点, 分支数没有增加. 因此把桥与割点混为一谈会在小图上立即出错.
\originalsection{不相交路与 Menger 定理}
\begin{theorem}[Menger 的边版本]\label{thm:mengeredge}
在有限无向图中, 两不同顶点 $s,t$ 之间两两边不相交路的最大条数, 等于使 $s,t$ 分离的最小边割大小.
\end{theorem}
\begin{proof}
每条 $s$--$t$ 路都必须穿过边割, 边不相交路不能共用该割中的边, 所以路数不超过割大小. 对反向, 把每条无向边换成两个反向弧, 容量各为1. 第\ref{ch:flow}章用独立的增广算法证明整数最大流最小割定理; 一个无向割的有向前向容量恰等于其边数. 取整数最大流, 若原边两个方向均有流量, 同时减去相同流量, 不改变任何顶点净流量. 剩余每条原边最多使用一个方向. 反复沿正流量从 $s$ 找到 $t$, 删去一单位路流; 若出现圈先消去圈流. 于是最大流分解为边不相交路, 数目等于最小割. 流定理的证明不依赖本章, 所以这里没有循环引用.
\end{proof}
\begin{theorem}[Menger 的顶点版本]\label{thm:mengervertex}
设 $s,t$ 为不相邻的不同顶点. 内部顶点两两不相交的 $s$--$t$ 路的最大条数, 等于可从 $V\setminus\{s,t\}$ 删除以分离 $s,t$ 的最小顶点数.
\end{theorem}
\begin{proof}
每条路必须经过所删顶点之一, 因而得到上界. 为证明下界, 把每个非端点 $v$ 拆为 $v_{\rm in},v_{\rm out}$, 加容量1的弧 $v_{\rm in}\to v_{\rm out}$. 原边按可行方向变成从出端到入端的弧, 容量取 $M=n+1$, 端点 $s,t$ 不拆. 不相邻条件使删除全部其他顶点能分离它们, 所以存在容量至多 $n-2$ 的割. 最小割不可能含容量 $M$ 的弧, 只能由某些容量1的顶点弧组成. 这些顶点在原图构成分离集, 反之每个分离集给出不大于其大小的割.

整数最大流分解为路后, 每条容量1的顶点弧最多被一路使用, 投影回原图就是内部顶点不相交路. 删除投影中的可能重复段不增加共用顶点. 因而最大路数等于最小顶点割.
\end{proof}
“不相邻”不能删去: 若 $st$ 本来是一条边, 删除其他顶点永远无法分开 $s,t$, 但内部不相交路数仍有限. 对 $n\ge k+1$, $1\le k\le n-1$, 全局 $k$ 连通性等价于任意两点之间有 $k$ 条内部顶点不相交路. 非相邻端点的正向由 Menger 立即得到. 对相邻的 $s,t$, 若集合 $S\subseteq V\setminus\{s,t\}$ 在 $G-st$ 中分离它们且 $|S|<k-1$, 令 $A$ 为 $(G-st)-S$ 中含 $s$ 的分支. 若 $A\setminus\{s\}\ne\varnothing$, 删除 $S\cup\{s\}$ 后该集合仍与 $t$ 分离, 得大小小于 $k$ 的顶点割. 若 $A=\{s\}$, 则 $d_G(s)\le|S|+1<k$, 也与 $k$ 连通图最小度至少 $k$ 矛盾. 故删边图的最小分离集至少有 $k-1$ 点, 用非相邻版本再加单边路 $st$. 反向, 删除少于 $k$ 个顶点时, 对任意两个未删点, 至多截断相应数量的不相交路, 至少一路保留, 所以余图连通.

\originalsection{块树}
Brooks 染色定理需要在图中选两顶点一起删除, 仍保留连通性. 仅知道“无割点”不直接保证任意两点都能删除, 因而需要描述有割点图怎样拼接.
\begin{definition}
连通图的块包括极大二连通子图和每条桥组成的二顶点子图; 单顶点图单独作为一个块. 二连通图有至少三个顶点, 连通且无割点. 块的内部顶点指在该块中、却不是原图割点的顶点.
\end{definition}
\begin{lemma}\label{lem:blocktree}
连通图的不同块至多共有一个顶点; 有共有顶点时它是割点. 以块和割点为两类节点, 当割点属于某块就连边, 得到的关联图为树. 若原图有割点, 该树至少有两个叶块; 每个叶块恰含一个割点, 删除其任一内部顶点仍使原图连通.
\end{lemma}
\begin{proof}
先说明二连通子图的粘合性质. 若两个二连通子图至少共有两点, 删任一点后仍至少有一个共有点, 两子图的剩余部分都连通, 所以并图无割点且连通. 它们若为极大块, 必属同一块. 桥不能与二连通子图共用两端, 否则桥在圈上. 故不同块至多共有一点.

若两块共点 $v$, 却有一条避免 $v$ 的路连接两块的其他顶点, 可取路径的首次、末次接触点, 并把路径与两块合并. 删除任一点后, 经过 $v$ 的连接与避开 $v$ 的路径至少保留一种, 各块内部也保持可连接; 若块是桥则其端点可沿这条新圈连接. 合并得到更大的无割点子图, 违反极大性. 所以共有点必为割点.

每条边属于某块: 非桥边在圈上, 从圈扩张到极大二连通子图即可. 原图路径可沿所属块转移, 关联图于是连通. 若关联图有圈, 其上每个块有两个不同的连接割点. 圈上块的并图删去任一点后, 各二连通块的剩余部分仍连通, 桥块至多失去一个端点. 圈上的不同连接点至多失去一个, 其余连接仍沿圈形成通路, 因而并图无割点, 会合成更大块, 矛盾. 因而关联图为树. 删去割点所得各分支都含它的邻点; 这些邻边不可能属于同一二连通块, 所以割点节点至少关联两个块, 所以叶节点只能是块. 有割点时树有至少两叶, 叶块只有一个连接割点. 删除内部顶点, 二连通叶块的余部仍连通; 若它是桥, 删除内部端点只去掉叶子. 其他块均不受影响.
\end{proof}
\begin{example}
两个三角形只共有一个顶点 $v$. 求它的块、割点、$\kappa$ 与 $\lambda$.
\end{example}
\begin{solution}
两个三角形为两个块, $v$ 为唯一割点, 块树是“块--$v$--块”三节点路. 因而 $\kappa=1$. 每条边在三角形上, 无桥, 故 $\lambda\ge2$; 顶点 $v$ 之外各点度2, 故 $\lambda=2$. 边冗余不能替代顶点冗余.
\end{solution}
\originalsection{习题与解答}
\begin{exercise}\label{ex:conn1}
求 $K_{a,b}$ 的顶点连通度和边连通度, 其中 $1\le a\le b$ 且 $a+b\ge3$.
\end{exercise}
\begin{exercise}\label{ex:conn2}
证明二连通图任意两条不同边都在同一个圈上.
\end{exercise}
习题\ref{ex:conn1}的解答.
\begin{solution}
删去小侧全部 $a$ 点后, 大侧有至少两点而不连通, 所以 $\kappa\le a$. 删去少于 $a$ 点后两侧仍非空, 完全二部图仍连通, 所以 $\kappa=a$. 又 $\kappa\le\lambda\le\delta=a$, 得 $\lambda=a$.
\end{solution}
习题\ref{ex:conn2}的解答.
\begin{solution}
把两条边分别细分, 加新顶点 $x,y$. 细分保持二连通性: 删除原顶点时, 被细分边的剩余段仍附着在另一端; 删除新顶点相当于删原边, 而二连通图无桥. 由 Menger 得两条内部不相交的 $x$--$y$ 路, 合起来构成经过 $x,y$ 的圈. 两新顶点度2, 圈必使用其两条关联边, 恢复原边后得到所求圈.
\end{solution}
